\documentclass[11pt]{article}
\usepackage[T1]{fontenc}
\usepackage[utf8]{inputenc}
\usepackage{amsmath,amssymb,amsthm,amsfonts}
\usepackage[margin=2.4cm]{geometry}
\usepackage{xcolor}
\usepackage{booktabs}
\usepackage{enumitem}
\usepackage{hyperref}
\hypersetup{colorlinks=true,linkcolor=blue,citecolor=blue,urlcolor=blue}

\newcommand{\R}{\mathbb{R}}
\newcommand{\C}{\mathbb{C}}
\newcommand{\Reop}{\operatorname{Re}}
\newcommand{\Imop}{\operatorname{Im}}
\newcommand{\norm}[1]{\left\|#1\right\|}
\newcommand{\code}[1]{\texttt{#1}}

\title{Notas internas sobre la parte num\'erica\\
\large (documento de trabajo; \emph{no} forma parte del paper)}
\author{Registro detallado de lo que se hizo, por qu\'e, y con qu\'e resultados}
\date{\today}

\begin{document}
\maketitle

\noindent
Este documento explica, sin omitir detalles, todo el trabajo num\'erico realizado para
la subsecci\'on ``The optimal error and the $c/\log n$ law'' (Sec.~4.3 del paper) y para
el Teorema~3.2. Est\'a pensado para que nosotros entendamos exactamente qu\'e se
comput\'o, qu\'e errores hubo, c\'omo se corrigieron, y cu\'ales son las limitaciones.

\tableofcontents

\section{Objetivo y notaci\'on}

Trabajamos en el espacio de Hardy $H^2(\C_\alpha)$, $\C_\alpha=\{\Reop s>\alpha\}$, con
producto interno de frontera
\[
\langle F,H\rangle_\alpha=\frac1{2\pi}\int_{-\infty}^{\infty}F(\alpha+it)\,\overline{H(\alpha+it)}\,dt,
\qquad \norm{F}_\alpha^2=\langle F,F\rangle_\alpha .
\]
Los objetos del paper son, para $k\ge2$,
\[
G_k(s)=\bigl(k^{-s}-k^{-1}\bigr)\frac{\zeta(s)}{s},\qquad E(s)=\frac1s .
\]
Definimos dos cantidades:
\begin{itemize}[leftmargin=*]
  \item el \textbf{error de M\"obius} $s_n(\alpha)=\norm{S_n-E}_\alpha^2$, con
        $S_n=\sum_{k=2}^n\mu(k)G_k$ (coeficientes \emph{fijos});
  \item la \textbf{distancia \'optima}
        $d_n(\alpha)^2=\inf_{c_2,\dots,c_n}\norm{\sum_{k=2}^n c_kG_k-E}_\alpha^2$
        (coeficientes \emph{libres}).
\end{itemize}
Siempre $d_n(\alpha)^2\le s_n(\alpha)$. Todo el estudio fino es en la l\'inea cr\'itica
$\alpha=1/2$. El valor de referencia conjeturado (Nyman--Beurling--B\'aez-Duarte,
Balazard--Saias) es
\[
\sum_{\rho}\frac1{|\rho|^2}=2+\gamma-\log(4\pi)\approx0.0461 .
\]

\section{La cantidad correcta: por qu\'e coeficientes \emph{reales}}
\label{sec:bug}

Este fue el punto m\'as delicado y donde hubo un \textbf{error en el primer prototipo}
que corregimos.

\paragraph{La sutileza.} La norma de $H^2$ es $\frac1{2\pi}\int_{-\infty}^{\infty}$.
Como $E$ y los $G_k$ tienen coeficientes de Dirichlet \emph{reales}, valen las
reflexiones $E(\overline s)=\overline{E(s)}$, $G_k(\overline s)=\overline{G_k(s)}$.
Por lo tanto, \textbf{solo si los coeficientes $c_k$ son reales}, la funci\'on
$F=\sum c_kG_k-E$ hereda $F(\overline s)=\overline{F(s)}$, su m\'odulo al cuadrado es
\emph{par} en $t$, y entonces
\[
\norm{F}_\alpha^2=\frac1{2\pi}\int_{-\infty}^{\infty}|F(\alpha+it)|^2\,dt
=\frac1{\pi}\int_{0}^{\infty}|F(\alpha+it)|^2\,dt .
\]

\paragraph{El error.} El primer prototipo (\code{proto\_dn.py}) montaba una matriz de
dise\~no $M$ con los $G_k$ muestreados solo en el semieje $t\ge0$ (con peso
$\sqrt{w/\pi}$) y resolv\'ia $\min_c\norm{Mc-y}^2$ por \code{numpy.linalg.lstsq}
\emph{dejando $c$ complejo}. Con $c$ complejo, $F$ pierde la simetr\'ia de reflexi\'on
y $\frac1\pi\int_0^\infty|F|^2 \ne \frac1{2\pi}\int_{-\infty}^{\infty}|F|^2$: el
\code{lstsq} explota la libertad compleja para ajustar mejor \emph{el semieje},
dando un residuo \textbf{menor que la distancia real} de $H^2$. Es decir,
\emph{subestimaba} $d_n^2$.

\paragraph{La formulaci\'on correcta.} Hay que minimizar sobre $c\in\R^{\,n-1}$. Con
$A=(\langle G_j,G_k\rangle_\alpha)$ y $b=(\langle E,G_k\rangle_\alpha)$, y usando que la
parte imaginaria de $A$ es antisim\'etrica (no contribuye para $c$ real),
\[
\norm{\textstyle\sum c_kG_k-E}_\alpha^2
=c^{\top}\tilde A\,c-2c^{\top}\tilde b+\norm{E}_\alpha^2,\qquad
\tilde A=\Reop A,\ \ \tilde b=\Reop b,
\]
de donde
\[
\boxed{\ d_n(\alpha)^2=\norm{E}_\alpha^2-\tilde b^{\top}\tilde A^{-1}\tilde b\ }.
\]
Como $\tilde A$ es real, sim\'etrica y definida positiva, la soluci\'on $c=\tilde
A^{-1}\tilde b$ es \emph{real} de todos modos: no perdemos nada por restringir a
$c$ reales; simplemente hab\'ia que calcular con $\Reop A$ y $\Reop b$ (equivalente a
$\frac1\pi\int_0^\infty\Reop[F\overline H]$).

\paragraph{Verificaci\'on num\'erica del error.} Comparamos, en la misma malla fina,
el \code{lstsq} complejo (mal) contra el problema con $c$ real (montando el dise\~no
real $[\,\Reop M;\,\Imop M\,]$, que coincide exactamente con $\Reop A,\Reop b$):
\[
\begin{array}{r|ccc}
n & \text{lstsq complejo (MAL)} & \text{$c$ real (correcto)} & \text{sub-estim.}\\
\hline
40  & 1.19119\times10^{-2} & 1.27319\times10^{-2} & -6.4\%\\
80  & 1.02292\times10^{-2} & 1.09950\times10^{-2} & -7.0\%\\
160 & 8.75673\times10^{-3} & 9.35455\times10^{-3} & -6.4\%\\
320 & 7.54823\times10^{-3} & 7.95974\times10^{-3} & -5.2\%\\
640 & 6.61234\times10^{-3} & 7.02153\times10^{-3} & -5.8\%\\
\end{array}
\]
Los valores del paper son los de la columna correcta. \textbf{Consecuencia feliz:} tras
la correcci\'on, $d_n^2\log n$ \emph{sube} y pasa a \textbf{oscilar sobre} $0.0461$ (antes
quedaba sistem\'aticamente por debajo); el acuerdo con la conjetura mejora.

\section{Los productos internos en forma cerrada parcial}
\label{sec:formulas}

Sobre la frontera, con $s=\alpha+it$, escribimos $G_k(\alpha+it)=a_k(t)\,\psi(t)$ donde
\[
a_k(t)=k^{-\alpha}e^{-it\log k}-k^{-1},\qquad \psi(t)=\frac{\zeta(\alpha+it)}{\alpha+it}.
\]
Sea el \emph{peso} $\rho(t)=|\psi(t)|^2=\dfrac{|\zeta(\alpha+it)|^2}{\alpha^2+t^2}$ (par,
real, integrable). Definimos la transformada
\[
W(x)=\frac1{2\pi}\int_{-\infty}^{\infty}\rho(t)\,e^{itx}\,dt\quad(\text{real, par}).
\]
Un c\'alculo directo (desarrollando $a_j\overline{a_k}$ en cuatro t\'erminos) da la
\textbf{matriz de Gram en forma cerrada}:
\[
A_{jk}=(jk)^{-\alpha}\,W(\log k-\log j)
-\frac{j^{-\alpha}}{k}W(\log j)
-\frac{k^{-\alpha}}{j}W(\log k)
+\frac{1}{jk}W(0).
\]
Para el vector $b$, con $g(t)=\dfrac{\overline{\psi(t)}}{\alpha+it}
=\dfrac{\overline{\zeta(\alpha+it)}}{\alpha^2+t^2}$ (hermitiana: $g(-t)=\overline{g(t)}$,
luego su transformada es \emph{real}) y
$V(x)=\frac1{2\pi}\int_{-\infty}^{\infty}g(t)e^{itx}dt$,
\[
b_k=k^{-\alpha}V(\log k)-k^{-1}V(0).
\]
Adem\'as $\norm{E}_\alpha^2=\frac1{2\pi}\int\frac{dt}{\alpha^2+t^2}=\frac1{2\alpha}$
(en $\alpha=1/2$ vale $1$; sirve de test de la cuadratura). \textbf{Todo $A$ y $b$ son
reales}, coherente con la Secci\'on~\ref{sec:bug}.

\section{M\'etodo 1: denso estable (el que produjo los n\'umeros del paper)}
\label{sec:denso}

Implementado en \code{dn\_run.py}. Es el m\'etodo \textbf{de referencia y el que se
report\'o}.
\begin{enumerate}[leftmargin=*]
  \item \textbf{Malla.} $\approx2.8\times10^4$ nodos en $t\ge0$: parte fina uniforme
        $t\in[0,80]$ con paso $0.01$, m\'as cola \emph{geom\'etrica} hasta
        $t=3\times10^5$ ($2\times10^4$ nodos). Pesos trapezoidales.
  \item \textbf{Ensamblado.} $M$ de columnas $G_k\sqrt{w/\pi}$, $y=E\sqrt{w/\pi}$;
        luego $\tilde A=\Reop(M^HM)$, $\tilde b=\Reop(M^Hy)$,
        $\norm{E}^2=\Reop(y^Hy)$.
  \item \textbf{Resoluci\'on.} $d_n^2=\norm{E}^2-\tilde b_{[:n]}^{\top}
        \tilde A_{[:n,:n]}^{-1}\tilde b_{[:n]}$ v\'ia \code{numpy.linalg.solve}
        sobre la submatriz principal (barato reusar el $\tilde A$ grande).
  \item \textbf{$s_n$.} Para $c=\mu(k)$: $s_n=\mu^{\top}\tilde A\mu-2\mu^{\top}\tilde b+\norm{E}^2$
        (los $\mu(k)$ ya son reales, aqu\'i \emph{no} hab\'ia bug).
\end{enumerate}

\paragraph{Validaciones hechas.}
\begin{itemize}[leftmargin=*]
  \item $\norm{E}_{1/2}^2=0.999999$ contra el exacto $1$ (6 d\'igitos) $\Rightarrow$
        la cuadratura integra bien.
  \item Coincidencia exacta entre \code{solve}$(\Reop A,\Reop b)$ y el dise\~no real
        $[\Reop M;\Imop M]$ v\'ia \code{lstsq} (misma respuesta, confirma la
        formulaci\'on de la Secci\'on~\ref{sec:bug}).
  \item Estabilidad de $d_n^2$ bajo refinamiento de malla (probado con mallas de
        $7\times10^3$ hasta $4.2\times10^4$ nodos): las curvas coinciden a precisi\'on
        de gr\'afico.
\end{itemize}

\section{M\'etodo 2: acelerado por FFT (explorado, \emph{no} usado en el paper)}
\label{sec:fft}

Implementado en \code{dn\_gram.py}. La idea era usar las f\'ormulas cerradas de la
Secci\'on~\ref{sec:formulas} para llegar a $n\sim10^4$--$10^5$ sin montar la $M$ gigante:
muestrear $\rho(t)$ y $g(t)$ en malla uniforme, obtener $W,V$ por FFT, e interpolar en
$x=\log(k/j)$ y $x=\log k$.

\paragraph{Un bug de signo (corregido).} La transformada usa $e^{+itx}$. Para $\rho$
\emph{real}, \code{np.fft.fft} (que lleva $e^{-i\cdots}$) da lo mismo en la parte real,
as\'i que $W$ sal\'ia bien. Pero para $g$ \emph{complejo}, el t\'ermino con $\Imop g$
cambia de signo: hab\'ia que usar $M\cdot\code{ifft}$ (que lleva $e^{+i\cdots}$).
Sin la correcci\'on, $V$ (y por tanto $b$) sal\'ia mal por un factor grande.
Tras corregir, $A$ y $b$ coincid\'ian con el denso.

\paragraph{Por qu\'e \emph{no} escala (raz\'on de fondo).} Dos muros:
\begin{enumerate}[leftmargin=*]
  \item \textbf{Cancelaci\'on catastr\'ofica.} $d_n^2=\norm{E}^2-\text{proy}$ con
        $\norm{E}^2\approx1$ y $\text{proy}\approx0.99$: restar dos n\'umeros casi
        iguales. Para obtener $d_n^2\sim10^{-3}$ a 2--3 cifras hace falta la
        proyecci\'on a $\sim6$ cifras.
  \item \textbf{Cola del peso.} $\rho(t)\sim\log t/t^2$; truncar en $T$ deja un residuo
        $\sim\log T/T$ en $W$ que, amplificado por la cancelaci\'on, contamina $d_n^2$.
\end{enumerate}
Medido (referencia = denso correcto):
\[
\begin{array}{ll|ccc}
dt & T & n{=}160 & n{=}640 & n{=}1280\\
\hline
0.02 & 2\times10^4 & +3.7\% & +5.9\% & +7.0\%\\
0.01 & 2\times10^4 & +3.7\% & +5.9\% & +7.0\%\\
0.01 & 4\times10^4 & +1.5\% & +2.4\% & +2.9\%\\
0.005& 4\times10^4 & +1.5\% & +2.4\% & +2.9\%\\
\end{array}
\]
Lectura: el paso $dt$ \emph{no} importa (la malla en $t$ ya basta); el error lo fija
\emph{$T$} (la cola). Duplicar $T$ redujo el error a la mitad. Para $<0.5\%$ har\'ia falta
$T\sim10^6$, es decir $\sim5\times10^7$ nodos: demasiado caro para la ganancia. Por eso
\textbf{descartamos el M\'etodo~2 para producir n\'umeros} y nos quedamos con el denso.

\section{Condicionamiento y por qu\'e $n\le5120$}
\label{sec:cond}

La matriz $\tilde A=M^HM$ tiene $\operatorname{cond}(\tilde A)=\operatorname{cond}(M)^2$;
los $G_k$ se vuelven casi linealmente dependientes al crecer $k$. N\'umeros medidos en
la malla fina:
\[
\begin{array}{r|cc}
n & \operatorname{cond}(M) & \operatorname{cond}(\tilde A)\\
\hline
80   & 1.84\times10^{2} & 3.40\times10^{4}\\
160  & 4.12\times10^{2} & 1.70\times10^{5}\\
320  & 8.88\times10^{2} & 7.88\times10^{5}\\
640  & 1.93\times10^{3} & 3.73\times10^{6}\\
1280 & 4.54\times10^{3} & 2.06\times10^{7}\\
2560 & 1.25\times10^{4} & 1.57\times10^{8}\\
\end{array}
\]
A $n=2560$, $\operatorname{cond}(\tilde A)\sim10^{8}$: el \code{solve} pierde $\sim8$
d\'igitos de los $\sim16$ disponibles, sobran para $d_n^2$ a 3--4 cifras. A
$n=5120$ (\'ultimo punto reportado) segu\'imos en zona estable; m\'as all\'a de eso
la combinaci\'on cancelaci\'on $+$ condicionamiento har\'ia falta multiprecisi\'on.
\emph{Adem\'as} la ganancia cient\'ifica de subir $n$ es marginal: la convergencia
$1/\log n$ es glacial ($1/\log n$ pasa de $0.117$ en $n{=}5120$ a $0.087$ en $n{=}10^5$).

\section{Resultados finales (corregidos)}
\label{sec:result}

Malla de $2.8\times10^4$ nodos, $\alpha=1/2$, hasta $n=5120$:
\[
\begin{array}{r|ccc}
n & d_n(1/2)^2 & s_n(1/2) & d_n^2\log n\\
\hline
10   & 2.38407\times10^{-2} & 1.09220\times10^{-1} & 0.05490\\
20   & 1.65317\times10^{-2} & 8.95797\times10^{-2} & 0.04952\\
40   & 1.27319\times10^{-2} & 1.44170\times10^{-1} & 0.04697\\
80   & 1.09950\times10^{-2} & 9.25434\times10^{-2} & 0.04818\\
160  & 9.35455\times10^{-3} & 6.84596\times10^{-2} & 0.04748\\
320  & 7.95974\times10^{-3} & 5.35622\times10^{-2} & 0.04591\\
480  & 7.43951\times10^{-3} & 7.57841\times10^{-2} & 0.04593\\
640  & 7.02153\times10^{-3} & 3.65723\times10^{-2} & 0.04537\\
960  & 6.61986\times10^{-3} & 7.94786\times10^{-2} & 0.04546\\
1280 & 6.36071\times10^{-3} & 3.47250\times10^{-2} & 0.04551\\
1920 & 6.12159\times10^{-3} & 4.17903\times10^{-2} & 0.04628\\
2560 & 5.93084\times10^{-3} & 4.31240\times10^{-2} & 0.04654\\
3200 & 5.58352\times10^{-3} & 9.45740\times10^{-2} & 0.04506\\
4096 & 5.38661\times10^{-3} & 1.22126\times10^{-1} & 0.04480\\
5120 & 5.28239\times10^{-3} & 3.85040\times10^{-2} & 0.04512\\
\end{array}
\]
Ajuste $d_n^2\sim C/\log n+D/(\log n)^2$ por regresi\'on en $1/\log n$:
$C=0.0407$, $D=0.0306$. El $C$ l\'ider queda un $\sim12\%$ por debajo de $0.0461$ porque
el rango de $n$ es corto (el t\'ermino $D/(\log n)^2$ absorbe parte); pero
$d_n^2\log n$ \emph{crudo} ya oscila en $0.045$--$0.047$, cruzando $0.0461$. Esa es la
historia honesta: acuerdo a pocos \% con la constante conjeturada, con la salvedad de
que fijar $C$ con precisi\'on est\'a fuera de alcance por la lentitud $1/\log n$.

\paragraph{Contraste $\mu(k)$ vs \'optimo.} $s_n$ se queda en $[3.5\times10^{-2},
1.4\times10^{-1}]$ oscilando sin decaer (es lo que predice B\'aez-Duarte: los
coeficientes fijos $\mu(k)$ \emph{divergen}); $d_n^2$ decae monot\'ono. El cociente
$s_n/d_n^2$ va de $\sim5$ a $\sim20$. Ese contraste es la figura central (Fig.~4 del
paper).

\section{Qu\'e se a\~nadi\'o al paper y c\'omo reproducir}

\paragraph{En \code{riemann.tex}.}
\begin{itemize}[leftmargin=*]
  \item \textbf{Subsecci\'on 3.3} ``A lower bound from the zeros of $\zeta$'':
        Teorema~3.2 (cota inferior por n\'ucleo reproductor,
        $\norm{E-\sum c_kG_k}_\alpha^2\ge 2(\Reop\rho-\alpha)/|\rho|^2$, y versi\'on
        multi-cero $d_n^2\ge v^{*}K^{-1}v$) y Corolario~3.2
        ($\tfrac12\le\Theta\le\alpha_c\le1$; $\alpha_c=\tfrac12\Rightarrow$ RH).
  \item \textbf{Subsecci\'on 4.3} ``The optimal error and the $c/\log n$ law'':
        m\'etodo, Figura~4 (contraste $+$ meseta) y Tabla~4 (valores corregidos).
\end{itemize}

\paragraph{Scripts (en la ra\'iz del repo).}
\begin{itemize}[leftmargin=*]
  \item \code{dn\_run.py} --- c\'omputo de referencia (M\'etodo~1). Genera
        \code{dn\_results.npz} y \code{dn\_results.csv}.
  \item \code{dn\_figure.py} --- produce \code{latex/fig\_dn.pdf} desde el
        \code{.npz}.
  \item \code{dn\_gram.py} --- M\'etodo~2 (FFT) con su validaci\'on; queda como
        respaldo, \emph{no} se us\'o para los n\'umeros.
  \item \code{proto\_dn.py} --- prototipo inicial (conten\'ia el bug de coef.
        complejos; se conserva solo como historia).
\end{itemize}
Dependencias: Python~3, \code{numpy}, \code{scipy}$\ge$1.17 (soporte complejo de
\code{scipy.special.zeta}), \code{matplotlib}. Reproducir: \code{python dn\_run.py \&\&
python dn\_figure.py}.

\paragraph{Puntos abiertos.}
\begin{itemize}[leftmargin=*]
  \item Barrido en $\alpha\in[1/2,1]$ para estimar $\alpha_c$ num\'ericamente
        (cierra el Corolario~3.2).
  \item Si se quisiera $n\gg5120$: multiprecisi\'on (\code{mpmath}) en el \code{solve},
        o correcci\'on anal\'itica de la cola de $\rho$ en el M\'etodo~2.
\end{itemize}

\end{document}
